\chapter{Conclusion and Outlook}

This report has provided an overview of the current state of research in the field of hand gesture recognition and highlighted some of its applications in sign language translation. Examining the relevant literature revealed a trend towards deep learning-based approaches, which have yielded significant improvements in accuracy and robustness. Another common theme in both chapters was the use of multimodal data to make up for the limitations of a single video feed. Supplementary depth maps and skeletal data can make up for ambiguity caused by lighting and occlusion in the image. 


Foregoing the gloss representation stage in sign translation has shown both promising results in accuracy as well as reduced the work required in creating suitable training data. For deep learning-based approaches, the availability of large, annotated datasets is crucial. The covered literature focuses on American and German sign language, which have a large number of speakers and thus a larger pool of data. Enabling translation of less popular sign languages, specific regional dialects and individual signing styles, requires a lower barrier for data collection and annotation. 

Automated sign recognition and production has the potential to improve accessibility online and in real life. In education, new tools could be developed to complement traditional teaching methods, which would allow for wider adoption of the language. It is, however, important to note that sign language is traditionally taught by non-hearing individuals in an effort to preserve the culture and history of a language that has been historically marginalized \cite{ASLTeachingEthics}. Besides the technical challenges, the social and cultural implications of who gets to build and control these systems will become increasingly important as the technology matures.